\documentclass[10pt,a4paper]{amsart}
\usepackage[margin=19mm]{geometry}
\usepackage{amsmath,amssymb,mathtools}
\usepackage[scr=rsfs,cal=euler]{mathalfa}
\usepackage{xfrac}
\usepackage[unicode,hidelinks,bookmarks=false]{hyperref}
\newcommand{\ie}{i.e.\ }
\newcommand{\cf}{cf.\ }
\newcommand{\resp}{respectively\ }
\newcommand{\Subsection}{\subsection}
\providecommand{\nobreakdash}{\nobreak\hbox{-}}
\setlength{\emergencystretch}{2em}
\multlinegap0pt
\newtheorem{lemma}{Lemma}
\title{Critical point localization and multiplicity results in Banach spaces via Nehari manifold technique}
\author{Radu Precup, Andrei Stan}
\date{Source: arXiv:2505.08678v2 (CC BY 4.0)}
\begin{document}
\maketitle

\noindent\emph{Source and licence.} Critical point localization and multiplicity results in Banach spaces via Nehari manifold technique. Version arXiv:2505.08678v2 (CC BY 4.0), distributed by arXiv under the Creative Commons Attribution 4.0 International licence, http://creativecommons.org/licenses/by/4.0/, as recorded in the arXiv licence metadata for this exact version and shown on the abstract page https://arxiv.org/abs/2505.08678v2. Full licence notice: \texttt{licenses/arxiv-2505.08678-CC-BY-4.0.txt}.\par\smallskip

\noindent\textit{Editorial scope.} Counted unit: the article's lemma copy(1) (source0.tex lines 425--436). The article's complete original proof of it is delivered with it (source0.tex lines 438--478, one \texttt{proof} environment of 41 lines). No mathematics is written, repaired or re-derived: the statement and the proof are the article's own lines, in the article's own notation, and no retained line is substituted or re-flowed.  Retained uncounted as the source-local context the proof needs: lines 421--423.  Nothing was omitted from any retained span, so the package carries no omission record.  No retained line contains a citation, so the wrapper carries no bibliography and no bibliography entry.  No retained line contains a figure environment, an image-inclusion command or drawing code, so no image is delivered and the package declares no asset dependency.  Editorial formatting only, in addition: an AMS \texttt{amsart} preamble that restates the article's own declarations and macros used by the retained lines, namely the environments \texttt{lemma} on separate counters, together with the standard packages those lines need.  The retained environments print in this document as Lemma 1 in order of appearance, whereas the article prints the same items with the numbering of its own full document.  The complete native source of the article as distributed in the arXiv e-print for this exact version is bundled unchanged as \texttt{sources/178057/source0.tex}.
\begin{equation}
\frac{r}{|u|}<s\left( u\right) <\frac{R}{|u|}.  \label{f5}
\end{equation}

\begin{lemma}
\label{lc copy(1)}Under assumption (H3), the Nehari manifold $\mathcal{N}%
_{r,R}$ has the representations:
\begin{equation*}
\mathcal{N}_{r,R}=\left\{ \,s(u)\,u:\ u\in K_{r,R}\,\right\} =\left\{ u\in
K_{r,R}:\ s\left( u\right) =1\right\} .
\end{equation*}%
Moreover, for every $u\in \mathcal{N}_{r,R},$ one has
\begin{equation}
r<\left\vert u\right\vert <R.  \label{f2}
\end{equation}
\end{lemma}

\begin{proof}
Denote
\begin{equation*}
A:=\left\{ \,s(u)\,u:\ u\in K_{r,R}\,\right\} ,\ \ \ B:=\left\{ u\in
K_{r,R}:\ s\left( u\right) =1\right\} .
\end{equation*}

(a) If $u\in \mathcal{N}_{r,R},$ one has $u\in K_{r,R}$ and $\langle
E^{\prime }(u),u\rangle =0,$ that is $\alpha _{u}^{\prime }\left( 1\right)
=0.$ Hence $s\left( u\right) =1$ and so $u=s\left( u\right) u\in A.$ Thus $%
\mathcal{N}_{r,R}\subset A.$

(b) Conversely, if $u\in A,$ then $u=s\left( v\right) v$ for some $v\in
K_{r,R},$ and from the definition of $s,$
\begin{equation*}
0=\alpha _{v}^{\prime }\left( s\left( v\right) \right) =\left\langle
E^{\prime }\left( s\left( v\right) v\right),v \right\rangle =\left\langle
E^{\prime }\left( u\right) ,v\right\rangle ,
\end{equation*}%
whence $\left\langle E^{\prime }\left( u\right) ,u\right\rangle =0,$ showing
that $u\in \mathcal{N}_{r,R}.$ Thus $A\subset \mathcal{N}_{r,R}.$

Therefore $\mathcal{N}_{r,R}=A.$

(c) The inclusion $\mathcal{N}_{r,R}\subset B$ is clear in view of (a).

(d) If $u\in B,$ then $s\left( u\right) =1,$ so $\alpha _{u}^{\prime }\left(
1\right) =0,$ i.e., $\left\langle E^{\prime }\left( u\right) ,u\right\rangle
=0,$ which shows that $u\in \mathcal{N}_{r,R}.$ Thus $B\subset \mathcal{N}%
_{r,R}.$

Therefore $\mathcal{N}_{r,R}=B.$

To prove (\ref{f2}), take any $u\in \mathcal{N}_{r,R}.$ Then $s\left(
u\right) =1,$ and in virtue of (\ref{f5}), one has
\begin{equation*}
\frac{r}{\left\vert u\right\vert }<s\left( u\right) =1<\frac{R}{\left\vert
u\right\vert },
\end{equation*}%
which yields (\ref{f2}).
\end{proof}

\end{document}
